% \section{Evo-Bench: A Benchmark for Measuring Harness-Evolving Capability}
\section{Evo-Bench}
\label{sec:benchmark}

Evo-Bench is a benchmark for evaluating large language models' \emph{harness-evolving capability}: 
the ability to conduct long-horizon, code-centric iterative improvement of an executable agent harness. Given a set of validation tasks, models are expected to diagnose failure patterns, formulate improvement hypotheses, and revise the harness implementation to achieve sustained performance gains.

\begin{figure}[!t]
    \centering
    \includegraphics[width=0.84\linewidth]{Figures/dataset_composition.pdf}
    \caption{Composition of Evo-Bench.  The left panel summarizes the
    three-domain hierarchy and source-benchmark allocation; the right panel
    compares the validation and evaluation splits.}
    \label{fig:dataset-composition}
\end{figure}


% the ability to conduct long-horizon, code-centric research that diagnoses failures, formulates improvement hypotheses, and revises an executable agent harness so that its gains transfer beyond visible development tasks. 
% to make it a robust, general harness


% As shown in Figure~\ref{fig:evlaution}, each Evo-Bench run couples a task-solving policy agent with an evolver. 
% Let $\pi$ denote the fixed policy model and $E$ the evolver model. 

% \textcolor{blue}{Starting from a shared seed harness $H_0$, the evolver iteratively inspects evidence accumulated on the visible validation tasks and modifies the current harness $H_t$ to produce $H_{t+1}$.
% Each revision is then formally evaluated on the same suite.}

% \textcolor{blue}{
% As shown in Figure~\ref{fig:evlaution}, each Evo-Bench run consists of a fixed policy model $\pi$ and an evolver $E$ that iteratively improves the policy harness.
% At evolution iteration $t$, the policy agent is $A_t^{\mathrm{task}}=(\pi,H_t)$, where the editable policy harness $H_t$ governs how $\pi$ interacts with target tasks.
% }
\paratitle{Task Formulation.}
As shown in Figure~\ref{fig:evaluation}, each Evo-Bench run consists of a fixed policy model $\pi$ and an evolver $E$ that iteratively improves the policy harness.
At evolution iteration $t$, the policy agent is $A_t^{\mathrm{task}}=(\pi,H_t)$, where the editable policy harness $H_t$ governs how $\pi$ interacts with target tasks.
The evolver is $A^{\mathrm{evo}}=(E,\mathcal{H}_{\mathrm{evo}})$ and operates through a separate, fixed evolve harness $\mathcal{H}_{\mathrm{evo}}$. 
Let $\mathcal{E}_t^{\mathrm{val}}:=\bigl((H_i,j_i^{\mathrm{val}},O_i^{\mathrm{val}})\bigr)_{i<t}$ denote the cumulative validation-side evidence available at iteration $t$, where $j_i^{\mathrm{val}}$ is the aggregate validation score and $O_i^{\mathrm{val}}$ contains the task-level outcomes, policy trajectories, and diagnostic feedback from evaluation $i$.
Starting from the common seed harness $H_0$, the evolver inspects the accumulated validation-side evidence, edits the current harness $H_t$, and requests a formal evaluation of the resulting revision on the visible validation suite at each iteration. 
The evolution process is constrained by a fixed resource budget $\mathbf{b}=(b^{\mathrm{iter}},b^{\mathrm{time}},b^{\mathrm{steps}})$, which bounds the number of iterations, wall-clock time, and evolver steps. When the run terminates, the final revision $H_T$ is frozen and evaluated on a disjoint held-out evaluation suite.
% Evolution proceeds under a fixed resource budget
% at each iteration $t$
% formal evaluations
% externally on the separated eval suite.


% \paratitle{Composition.}
% \label{sec:benchmark-composition}
% Evo-Bench consists of a 160-task visible validation suite $\mathcal{D}_{\mathrm{val}}$ and a separated 448-task eval suite $\mathcal{D}_{\mathrm{eval}}$. The evolver can request formal evaluations only on $\mathcal{D}_{\mathrm{val}}$; the eval suite is reserved for external evaluation after the harness is frozen. Both suites draw from five source benchmarks: BrowseComp~\citep{browsecomp} and HLE~\citep{hle} form the search domain, GDPval~\citep{gdpval} and APEX-Agents~\citep{apexagents} form the office domain, and Claw-Eval~\citep{claweval} forms the general domain. The validation suite contains 32 tasks from each source. The eval suite contains 128 BrowseComp tasks, 128 HLE tasks, 64 GDPval tasks, 64 APEX-Agents tasks, and 64 Claw-Eval tasks.


\paratitle{Task Domains and Composition.}
% \textcolor{blue}{
\label{sec:benchmark-composition}
Evo-Bench covers three representative agent domains: search, office, and general agent tasks, evaluating harness evolution across diverse interaction patterns and tool-use requirements.
To instantiate these domains, we select five  established and challenging benchmarks: BrowseComp~\citep{browsecomp} and HLE~\citep{hle} for search, GDPval~\citep{gdpval} and APEX-Agents~\citep{apexagents} for office, and Claw-Eval~\citep{claweval} for general agent tasks. 
Based on these benchmarks, we construct a 160-task visible validation suite $\mathcal{D}_{\mathrm{val}}$ and a disjoint 448-task evaluation suite $\mathcal{D}_{\mathrm{eval}}$. 
Each source benchmark contributes 32 validation tasks, while the evaluation suite contains 128 BrowseComp, 128 HLE, 64 GDPval, 64 APEX-Agents, and 64 Claw-Eval tasks.
During evolution, the evolver can request validation evaluations only on $\mathcal{D}_{\mathrm{val}}$, while $\mathcal{D}_{\mathrm{eval}}$ is reserved for final evaluation after the harness is frozen.
Figure~\ref{fig:dataset-composition} summarizes the domain hierarchy and the
per-source allocation across the two splits.

% Evo-Bench consists of a 160-task visible validation suite $\mathcal{D}_{\mathrm{val}}$ and a separated 448-task eval suite $\mathcal{D}_{\mathrm{eval}}$.
% The evolver can request validation evaluations on $\mathcal{D}_{\mathrm{val}}$; the evaluation suite is reserved for external evaluation after the harness is frozen.
% Both suites draw from five source benchmarks: BrowseComp~\citep{browsecomp} and HLE~\citep{hle} form the search domain, GDPval~\citep{gdpval} and APEX-Agents~\citep{apexagents} form the office domain, and Claw-Eval~\citep{claweval} forms the general domain. The validation suite contains 32 tasks from each source. The eval suite contains 128 BrowseComp tasks, 128 HLE tasks, 64 GDPval tasks, 64 APEX-Agents tasks, and 64 Claw-Eval tasks.
% }

\paratitle{Evolve Harness.}
\label{sec:evolve-harness}
% we integrate specialized skills tailored specifically for the evolution process, 
The evolve harness operates as a fixed agent loop explicitly designed to support long-horizon autonomous work. Its core components, including the system prompt, execution tools, and context manager, are inspired by the design of Claude Code. Building upon this foundation, we integrate specialized skills for the evolution process, such as trajectory analysis and experiment tracking. We implement this custom harness natively within Evo-Bench to enable a controlled and transparent evaluation protocol.
% rather than relying on external harness implementations like Claude Code or Codex. This design enables a controlled and transparent evaluation protocol.

\paratitle{Policy Harness.}
\label{sec:policy-harness}
The policy harness is the evolving artifact optimized by the evolver. Its initial version $H_0$ is a minimal CodeAct~\cite{wang2024executable} loop with only a shell-execution tool and a final-answer completion tool. During validation evaluation on $\mathcal{D}_{\mathrm{val}}$, the fixed policy model executes tasks in a separate sandbox through the current harness $H_t$.
The evolver aims to transform this initial harness into a more capable, generalizable harness that boosts performance across all three domains.

\begin{figure}[!t]
    \centering
    \includegraphics[width=0.92\linewidth]{Figures/harness_evaluation.pdf}
    \caption{Overview of the Evo-Bench evaluation pipeline.}
    \label{fig:evaluation}
\end{figure}

\paratitle{Evaluation Metrics.}
\label{sec:benchmar_metrics}
We evaluate harness-evolving capability from two complementary perspectives: final generalization performance and evolutionary progress.
For a given suite $\mathcal{D}$, let $\mathcal{S}(\pi, H; \mathcal{D})$ denote the aggregate score achieved by the policy model with harness $H$, computed from the source benchmarks' native scorer~(defined in Table~\ref{tab:task-selection-stats}) and aggregated across target domains. 

To measure the final performance of the evolved harness, we define \textbf{Overall Score} on the unseen evaluation suite:
\[
\mathrm{Overall}(E)
:=\mathcal{S}(\pi,H_T;\mathcal{D}_{\mathrm{eval}}).
\]

To measure the evolution progress, we define \textbf{Anytime Validation Score} as the average best-so-far validation performance over the evolution budget. At iteration $t$, the best validation score achieved up to and including iteration $t$ is:
\[
S_t^*=\max_{i\leq t}\mathcal{S}(\pi,H_i;\mathcal{D}_{\mathrm{val}}).
\]

The Anytime Validation Score is then defined as:
\[
\mathrm{AnytimeVal}(E)
:=\frac{1}{b^{\mathrm{iter}}}
\sum_{t=1}^{b^{\mathrm{iter}}}S_t^* .
\]

If a run terminates early, its final best-so-far validation score is carried forward for the remaining budget.

\begin{table}[htbp]
\centering
% \small
\caption{Statistics of the harness-guided task selection process. Tasks with non-positive sensitivity ($\mathrm{Sens}(x) \le 0$) are filtered out prior to sensitivity-aware stratified selection. \textbf{Metrics}: $\mathrm{Sens}(x)$ denotes harness sensitivity (Pearson correlation); $\mathrm{Perf}(x)$ denotes average performance across the auxiliary harness set $\mathcal{H}_{\mathrm{aux}}$. ``LJ'' represents LLM-as-a-Judge.}
\label{tab:task-selection-stats}
\resizebox{\textwidth}{!}{%
\begin{tabular}{l ll cc cc cc cc}
\toprule
& \multicolumn{2}{c}{\textbf{Evaluation Protocol}} & \multicolumn{2}{c}{\textbf{Candidate Filtering}} & \multicolumn{2}{c}{\textbf{Final Allocation}} & \multicolumn{2}{c}{\textbf{$\mathrm{Sens}(x)$}} & \multicolumn{2}{c}{\textbf{$\mathrm{Perf}(x)$}} \\
\cmidrule(lr){2-3} \cmidrule(lr){4-5} \cmidrule(lr){6-7} \cmidrule(lr){8-9} \cmidrule(lr){10-11}
{Source Dataset} & Evaluator & Metric & Candidates & $\mathrm{Sens}(x) \le 0$  & Validation & Evaluation & Mean & Median & All & Selected \\
\midrule
APEX-Agents & Rubric LJ & Pass@1 & 421 & 133 & 32 & 64  & 0.377 & 0.423 & 0.389 & 0.224 \\
BrowseComp  & LJ        & Pass@1 & 768 & 47  & 32 & 128 & 0.256 & 0.236 & 0.444 & 0.271 \\
Claw-Eval   & LJ+Rule   & Pass\^{}$3$ & 157 & 60  & 32 & 64  & 0.386 & 0.317 & 0.839 & 0.747 \\
GDPval      & Rubric LJ & Mean    & 215 & 95  & 32 & 64  & 0.368 & 0.336 & 0.797 & 0.495 \\
HLE         & LJ        & Pass@1 & 768 & 62  & 32 & 128 & 0.313 & 0.300 & 0.333 & 0.251 \\
\bottomrule
\end{tabular}
}
\end{table}
